\section{Methodology}
\label{sec:method}

We propose \cubic, a curvature-driven dynamic caching framework for
Flow Matching inference. The key idea is to connect local trajectory
geometry with runtime caching decisions: regions with locally smooth
trajectories can tolerate more aggressive cache reuse, whereas regions
with larger curvature should be refreshed more frequently. To make this
principle practical, \cubic first models the cache-induced state error and
derives a curvature-aware interval allocation rule. It then determines an
appropriate cache granularity in token space and estimates local
curvature online using a zero-overhead endpoint prediction drift proxy.
\subsection{Overview}
\label{sec:method_overview}

Recall from Sec.~\ref{sec:motiv_theory} that caching intermediate features for $\tau$ steps is mathematically equivalent to extending the effective local integration step size by a factor of $\tau$. \cubic exploits the spatial heterogeneity of permissible step sizes by alternating between two types of steps during the ODE solving process:
\begin{itemize}
    \item \textbf{Anchor steps}: full forward passes that refresh the exact velocity prediction for all tokens, update the curvature proxy, and re-evaluate the caching policy.
    \item \textbf{Selective steps}: sparse forward passes that recompute only a subset of tokens (the active tokens), while the remaining tokens (the frozen tokens) reuse cached intermediate states and outputs.
\end{itemize}
This anchor-selective loop is governed by three consecutive operations at each anchor step: measuring the curvature proxy, partitioning the token space, and optimizing block-specific update intervals.
% \subsection{Overview}
% \label{sec:method_overview}

% We propose \cubic, a curvature-driven dynamic caching framework for Flow
% Matching inference. Recall from Sec.~\ref{sec:motiv_theory} that reusing
% cached intermediate features for $\tau$ steps effectively enlarges the
% local integration step size by a factor of $\tau$. Since different spatial
% regions exhibit different trajectory curvatures, they also have different
% tolerances to cache reuse: locally smooth regions can be reused more
% aggressively, whereas high-curvature regions should be refreshed more
% frequently.

% \cubic realizes this principle through an anchor-selective sampling loop.
% At \textbf{anchor steps}, a full forward pass refreshes velocity
% predictions, updates the curvature proxy, and re-evaluates the caching
% policy. At \textbf{selective steps}, only active tokens are recomputed,
% while frozen tokens reuse cached intermediate states and outputs. Each
% anchor step then applies three operations: curvature proxy estimation,
% cache granularity selection, and curvature-aware interval optimization.
\subsection{Curvature-induced Cache Error and Interval Optimization}
\label{sec:error_optimization}

\paragraph{\textbf{Error Modeling.}}
During Flow Matching inference, the sampler advances the current state
according to the velocity field predicted by the model. Consider a latent
position $i$. Suppose a full computation is performed at time $t_0$,
yielding the cached velocity
\begin{equation}
    v_0^{(i)}
    =
    v_\theta^{(i)}(x(t_0), t_0).
    \label{eq:cached_velocity}
\end{equation}
If this cached result is reused for the following $\tau_i$ solver steps,
the local update can be viewed as a zero-order hold approximation of the
velocity field. Comparing this approximate trajectory with the exact
trajectory, we define the cache-induced state error as
\begin{equation}
    \epsilon_i(t)
    =
    \tilde{x}^{(i)}(t)
    -
    x^{(i)}(t).
    \label{eq:cache_state_error}
\end{equation}

Within a short reuse window, if the local trajectory curvature remains
approximately stable, the variation of the velocity field along the
trajectory is controlled by the local curvature $\kappa_i$. The
zero-order hold approximation then yields the following error scaling (See details in ~\ref{app:error_modeling}):
\begin{equation}
    \left\|
    \epsilon_i(t_0 + \tau_i \Delta t)
    \right\|
    \lesssim
    \frac{1}{2}
    \kappa_i
    \tau_i^2
    \Delta t^2.
    \label{eq:cache_error_bound}
\end{equation}
This relation shows that the cache-induced error grows with both local
curvature and reuse length. Hence, smoother latent positions can tolerate
longer reuse, while high-curvature positions should be refreshed more
frequently.

\paragraph{\textbf{Interval Optimization.}}
We next convert the above error model into an interval allocation rule
over latent positions. Let $c_i$ denote the computational cost of
recomputing latent position $i$. Since position $i$ is recomputed once
every $\tau_i$ steps, its average per-step computation cost is
$c_i / \tau_i$. Under a target computation budget $C$, we determine the
reuse intervals by solving
\begin{equation}
    \min_{\{\tau_i\}_{i=1}^{N}}
    \sum_{i=1}^{N}
     \left\|\alpha_i \kappa_i \tau_i^2\right\|^2,
    \quad
    \mathrm{s.t.}
    \quad
    \sum_{i=1}^{N}
    \frac{c_i}{\tau_i}
    =
    C .
    \label{eq:interval_optimization_position}
\end{equation}

The resulting (~\ref{app:optimal_interval}) closed-form optimum satisfies $\tau_b^\star \propto \kappa_b^{-\alpha/(\beta+1)}.$
Using the theoretical default $(\alpha,\beta)=(2,4)$ yields\todo{check.}
\begin{equation}
    \tau_b^\star \propto \kappa_b^{-2/5}.
\end{equation}
Thus, \cubic naturally assigns shorter intervals to high-curvature blocks and longer intervals to low-curvature ones.

\subsection{Cache Granularity Design}
\label{sec:partitioner}

The previous subsection gives an ideal position-wise allocation rule:
latent positions with larger curvature should be refreshed more
frequently. In practice, however,
assigning an independent cache policy to every token would introduce
substantial scheduling and indexing overhead. Therefore, \cubic first
determines an appropriate cache granularity in token space and then
applies the interval allocation rule at the block level.

Specifically, \cubic defines block partitions directly in token space. At each anchor step, it selects a global block size $\mathbf{b}=(b_f,b_h,b_w)$ from a candidate set and tiles the entire token grid uniformly with that size. Thus, the partition is adaptive in \emph{granularity} but uniform in \emph{tiling}.

For any candidate block size $\mathbf{b}$ that induces a partition $\mathcal{P}(\mathbf{b})=\{B_j(\mathbf{b})\}$, we measure intra-block consistency by the coefficient of variation (CV). Unlike absolute variance, CV captures relative fluctuations, making it invariant to the absolute scale of curvature:
\begin{equation}
    \mathrm{cv}_j(\mathbf{b}) =
    \frac{\mathrm{Std}(\{\kappa_i\}_{i \in B_j(\mathbf{b})})}
    {\mathrm{Mean}(\{\kappa_i\}_{i \in B_j(\mathbf{b})}) + \varepsilon}.
\end{equation}
The selection principle is simple: \textbf{choose the coarsest block size whose induced blocks remain internally curvature-consistent}. In practice, this criterion is approximated using block-wise CV statistics such as mean-CV and p90-CV. If a candidate already satisfies the consistency threshold, \cubic prefers the largest block size; otherwise, it falls back to the candidate with the best consistency statistics.

Once the partition is fixed, we aggregate token-wise curvature into a representative block curvature. Rather than using the block mean, we adopt a conservative high quantile to prevent a few high-curvature tokens from being drowned out by the block average:
\begin{equation}
    \kappa_b = \mathrm{Quantile}_{q}\big(\{\kappa_i\}_{i \in B_b}\big), \qquad q \approx 0.9.
\end{equation}

\subsection{Zero-overhead Curvature Proxy}
\label{sec:curvature_proxy}

The block-wise caching policy requires estimating the local curvature
$\kappa_b$ during inference. Directly computing trajectory curvature,
however, would require Jacobian or higher-order derivative information,
which is impractical for large-scale Flow Matching models. We therefore
introduce a zero-overhead block-wise curvature proxy based on endpoint
prediction drift.

For a block $b$, the endpoint prediction at time $t$ is defined as
\begin{equation}
    \hat{x}_{1}^{(b)}(t)
    =
    x^{(b)}(t)
    +
    (1 - t)
    v_{\theta}^{(b)}(x(t), t),
    \label{eq:block_endpoint_prediction}
\end{equation}
where $x^{(b)}(t)$ and $v_{\theta}^{(b)}(x(t),t)$ denote the latent state
and predicted velocity restricted to block $b$, respectively. If the
local trajectory of this block is nearly linear, its endpoint prediction
should remain stable across neighboring anchor steps. In contrast, a
curved block trajectory leads to a noticeable drift in the predicted
endpoint.

Let $t_{a_{m-1}}$ and $t_{a_m}$ denote two consecutive anchor steps. We
define the block-wise endpoint prediction drift as
\begin{equation}
    \Delta \hat{x}_{1}^{(b,m)}
    =
    \hat{x}_{1}^{(b)}(t_{a_m})
    -
    \hat{x}_{1}^{(b)}(t_{a_{m-1}}).
    \label{eq:block_endpoint_drift}
\end{equation}
This drift approximates the local token-wise trajectory curvature (~\ref{app:curvature_proxy}):
\begin{equation}
    \left\| \Delta \hat{x}_{1}^{(b,m)}\right\|
    \approx
    \kappa_i(t_{a_m})\,(1-t_{a_m})\,\Delta t_{a_m},
\end{equation}
where $\Delta t_{a_m} = |t_{a_m} - t_{a_{m-1}}|$ is the interval between two anchor observations. Hence, the normalized quantity $\| \Delta \hat{x}_{1}^{(b,m)}\| / \left[(1-t_{a_m})\Delta t_{a_m}\right]$ serves as a zero-overhead proxy for curvature. This proxy only uses quantities already available at anchor steps, namely
the current latent state and the model-predicted velocity. Therefore, it
introduces no additional model forward passes. The estimated block-wise
curvature $\hat{\kappa}_b^{(m)}$ is then used to instantiate the reuse
interval $\tau_b$ according to the curvature-aware allocation rule.

\subsection{Selective Token DiT Forward}
\label{sec:runtime}

Given the partition $\mathcal{P}$ and the block intervals $\{\tau_b\}$, the active blocks at selective step $k$ are defined as
\begin{equation}
    \mathcal{A}_k = \{ b \in \mathcal{P} \mid s_k \bmod \tau_b = 0 \},
\end{equation}
where $s_k$ is the local step index within the selective phase. This induces the active and frozen token sets
\begin{equation}
    \Omega_k^{\mathrm{act}} = \bigcup_{b \in \mathcal{A}_k} B_b,
    \qquad
    \Omega_k^{\mathrm{frz}} = \Omega \setminus \Omega_k^{\mathrm{act}}.
\end{equation}

At each Transformer layer, \cubic recomputes the Query, Key, Value, and FFN only for active tokens. Frozen tokens reuse cached per-layer states and final outputs. For active tokens, attention is evaluated against the full sequence by combining newly computed Key/Value tensors for active tokens with cached Key/Value tensors for frozen tokens:
\begin{equation}
    Y_\ell^{\mathrm{act}} =
    \mathrm{Attn}\!\left(
        Q_\ell^{\mathrm{act}},
        K_\ell^{\mathrm{full}},
        V_\ell^{\mathrm{full}}
    \right).
\end{equation}
The final output of the selective step is then formed by merging newly computed active-token outputs with cached frozen-token outputs:
\begin{equation}
    v_k(\Omega) =
    \mathrm{Merge}\!\left(
        v_k^{\mathrm{act}}(\Omega_k^{\mathrm{act}}),
        v_{k-1,\mathrm{cache}}(\Omega_k^{\mathrm{frz}})
    \right).
\end{equation}
This selective forward engine is the component that actually yields computational savings: active tokens traverse the full DiT computation path, whereas frozen tokens participate only through cache reuse.

To ensure robustness, \cubic executes selective forward only when the required caches are already initialized and shape-consistent. Otherwise, the runtime automatically falls back to a full forward pass to refresh the cache state.

\subsection{End-to-End Inference Procedure}
\label{sec:end2end}

Putting the components together, \cubic integrates seamlessly into the standard sampling pipeline with minimal structural changes. The inference process is stabilized by standard bootstrap and termination mechanisms: it initiates with a few full forward passes to establish reliable curvature observations and cache states, enters the main anchor-selective loop dictated by the target speedup $\rho$, and concludes with full forward passes to preserve final-generation fidelity. 

This pipeline turns local trajectory geometry into a concrete execution policy: the monitor estimates curvature, the partitioner designs granularity, the optimizer assigns refresh intervals, and the runtime engine realizes those decisions as selective token computation.